\chapter*{Preface}
\addcontentsline{toc}{chapter}{Preface}

This volume studies the growth of the second moment for the Stratonovich
hyperbolic Anderson model with temporally correlated Gaussian noise. The main
question has two parts: what is the correct scale of growth, and what
variational quantity determines the leading constant? Answering these questions
requires us to keep track of the solution itself, its approximation by smooth
noises, and the relation between deterministic wave propagation and Gaussian
changes of measure.

The intended reader has a background in graduate probability and analysis.
Familiarity with Gaussian processes, Fourier analysis, Sobolev spaces and
bounded operators will be useful. The later chapters also use spectral theory
and Young integration. These subjects enter for specific purposes: covariance
estimates construct the random forcing, phase-space evolution retains both
displacement and velocity, and variational estimates compare amplification
with its Gaussian cost. Definitions and the required versions of imported
results are given where they enter the argument.

The volume can be read in several ways. The separate reading edition gives a
short account of the model, main statements and proof strategy, with links to
the detailed arguments here. Within this book, Chapter~\ref{S:book-setting}
introduces the problem and states the main results. Readers who want the
independent upper bound can proceed through Chapters~\ref{S:G0}
and~\ref{S:book-independent-upper}. The sharp-constant argument occupies
Chapters~\ref{S:book-gaussian-foundations}--\ref{S:book-constant}; the guide
that follows explains its dependencies. The opening of each chapter identifies
its purpose and the earlier material needed to enter it.

The appendices have a separate role. They preserve explicit low-order
calculations, shorter estimates under additional assumptions, covariance
comparisons and an obstruction to a proposed clock representation. These
calculations explain which features a successful argument must retain. They
also offer concrete entry points for readers who prefer examples before the
full construction. Their hypotheses and limitations remain visible, and their
placement does not make them prerequisites for every main theorem.

This is a working research manuscript. The exposition distinguishes written
results, imported statements, conditional consequences and open questions.
In particular, the future-work discussion separates higher-moment conjectures
from the exploratory question of a scaling limit for logarithmic wavefields.
The notation list, linked statements, navigation diagrams and subject index
are intended to help readers move between these parts without losing the
assumptions under which an argument is made.
